\documentclass[12pt]{article}
\usepackage[T1]{fontenc}
\usepackage[utf8]{inputenc}
\usepackage{lmodern}
\usepackage{textcomp}
\usepackage{enumitem}
\usepackage{geometry}
\geometry{margin=1in}
\begin{document}
\begin{center}
Answer Key - Science - K-12 Level Assessment 21 \
\small (Answers are ordered exactly as in the question paper.)
\end{center}

\section*{Multiple Choice}
\begin{enumerate}[label=\arabic*.]
\item \textbf{Answer:} D
\\ \textit{Options:} A) the distance of the planets from the sun; B) the depths of the major oceans on Earth; C) the amount of rainfall each day for a month; D) the percent of various materials in solid waste
\\ \textit{Note:} A circle graph, or pie chart, is ideal for displaying the percent of various materials in solid waste because it visually represents proportional relationships, allowing for an easy comparison of how each material contributes to the whole.
\item \textbf{Answer:} B
\\ \textit{Options:} A) It reduces the amount of carbon dioxide production.; B) It reduces the production of oxygen.; C) It decreases the greenhouse effect.; D) It decreases pollutants in the air.
\\ \textit{Note:} Cutting and burning forests eliminates trees, which are essential for photosynthesis, thereby significantly reducing the overall production of oxygen in the atmosphere.
\item \textbf{Answer:} A
\\ \textit{Options:} A) angle of tilt of the table; B) acidity of the water; C) amount of sediment eroded; D) rate that sediment is eroded
\\ \textit{Note:} The angle of tilt of the table should be constant because it directly affects the flow rate of water and the resulting erosion, making it essential to isolate the variable of water acidity to accurately measure its influence on erosion rates.
\item \textbf{Answer:} C
\\ \textit{Options:} A) one that eats meat; B) one that uses plants for shelter; C) one that lives under the ground; D) one that builds nests on the ground
\\ \textit{Note:} Animals that live underground are better equipped to survive and compete for resources after a fire because they are protected from the immediate destruction of the fire and can access food sources that might not be available to surface-dwelling animals.
\item \textbf{Answer:} B
\\ \textit{Options:} A) flexibility; B) reflectiveness; C) temperature; D) volume
\\ \textit{Note:} The correct answer is reflectiveness because the calm surface of the pond allows light to bounce off it, creating a clear reflection of the boy's face.
\end{enumerate}

\section*{True / False}
\begin{enumerate}[label=\arabic*.]
\setcounter{enumi}{5}
\item \textbf{Answer:} False
\\ \textit{Options:} A) True; B) False
\\ \textit{Note:} Mushrooms are fungi and obtain their nutrients through decomposition and absorption rather than by using sunlight and photosynthesis like plants do.
\item \textbf{Answer:} False
\\ \textit{Options:} A) True; B) False
\\ \textit{Note:} The statement is false because a gas changes directly into a solid through a process called deposition, which releases heat energy rather than absorbing it.
\item \textbf{Answer:} True
\\ \textit{Options:} A) True; B) False
\\ \textit{Note:} The foam ball model needs to revolve around the central ball to demonstrate how the moon moves in relation to the Earth and the Sun during a solar eclipse, accurately depicting the alignment necessary for the event to occur.
\item \textbf{Answer:} True
\\ \textit{Options:} A) True; B) False
\\ \textit{Note:} Honeybees are essential pollinators that transfer pollen from one flower to another, facilitating fertilization and enabling a wide variety of plants to reproduce and produce fruits and seeds.
\item \textbf{Answer:} False
\\ \textit{Options:} A) True; B) False
\\ \textit{Note:} Gravity on Earth is primarily caused by the mass of the Earth itself, which creates a gravitational pull, rather than the Earth's rotation around its axis.
\end{enumerate}

\section*{Long Form Response}
\begin{enumerate}[label=\arabic*.]
\setcounter{enumi}{10}
\item \textbf{Answer:} The smallest unit of copper that retains its elemental properties is an atom. A copper atom consists of protons, neutrons, and electrons, which are arranged in a specific way that gives copper its characteristic properties, such as conductivity and malleability. The structure of copper atoms allows them to easily lose one electron, making copper an excellent conductor of electricity. This property is crucial for its applications in electrical wiring, where efficient transmission of electrical current is essential. Thus, the atomic structure of copper directly influences its effectiveness in various technological applications.
\\ \textit{Note:} correct because it identifies the copper atom as the smallest unit that maintains elemental properties and explains how its electron configuration facilitates high conductivity and malleability, essential for its use in electrical wiring.
\item \textbf{Answer:} Thirst and sweating are crucial physiological responses during exercise that signal the need for hydration. When we exercise, our bodies generate heat, leading to an increase in body temperature. To cool down, we sweat, which results in the loss of fluids. As our body loses water, the brain triggers the sensation of thirst, prompting us to drink and replenish lost fluids. This process not only helps maintain optimal body function but also prevents dehydration, which can negatively impact performance and health. Recognizing these signals is essential for athletes and anyone engaging in physical activity to stay hydrated and perform at their best.
\\ \textit{Note:} correct because it accurately describes the physiological mechanisms of sweating and thirst, which work together to regulate body temperature and signal the need for fluid intake to prevent dehydration during exercise.
\end{enumerate}

\vfill
\noindent\textit{End of Answer Key}
\end{document}